% Mobile and Embedded Computing, Lab 2. Compile: latexmk -xelatex lab02.tex
\documentclass[10pt,aspectratio=169,t]{beamer}
\usetheme[lab]{upbminimal}

\course{Mobile and Embedded Computing}
\decklabel{Lab 2}
\title{A sensor reading over MQTT}
\subtitle{I2C in TinyGo, MQTT in the middle, a live value in Flutter}
\author{Dinu-Ștefan Rusu}
\institute{FILS, Universitatea Politehnica București}
\date{Week 4}

\begin{document}

\titleframe

\begin{frame}[eyebrow=Today]{What this lab is for}
  \vskip 2mm
  \begin{grid}[2]
    \cell[\faMicrochip]{Read a sensor and publish it}{Drive an I2C sensor from TinyGo with a driver from the TinyGo drivers repository, then publish the reading as JSON to your Mosquitto broker every two seconds.}
    \cell[\faMobile]{Show it live in the app}{Subscribe with mqtt\_client, push every message into a BLoC, render it in a BlocBuilder, and survive the broker or the board going away.}
  \end{grid}
  \vskip 6mm
  \taskmeta{Time}{2 hours}{Submit}{Push to your team repo, branch \code{lab-2}, before next lab}
\end{frame}

\begin{frame}[embedded,eyebrow=Setup]{Three set-ups, one contract}
  \vskip 1mm
  \begin{grid}[2]
    \cell[\faWifi]{ESP32-C3 or ESP32-S3}{Wi-Fi comes from the espradio package. The board publishes to the broker itself.}
    \cell[\faIcon{plug}]{Raspberry Pi Pico}{No Wi-Fi on the plain Pico. It prints the reading, and a laptop bridge forwards it.}
    \cell[\faIcon{desktop}]{Wokwi simulator}{No physical sensor, no route to your laptop. Use a simulated part or a counter, and a public test broker.}
    \cell[\faIcon{file-signature}]{The contract}{All three paths give the broker the same topic, the same JSON, one message every two seconds.}
  \end{grid}
  \vskip 2mm
  \muted{The broker is the Mosquitto you installed in Lab 1. Start it before anything else.}
\end{frame}

\begin{frame}[fragile,eyebrow=Setup]{Where your code goes}
  \begin{steps}
    \item Branch the team repository and add the two dependencies:
\begin{shell}
git checkout -b lab-2
cd firmware && go get tinygo.org/x/drivers
cd ../app   && flutter pub add mqtt_client
\end{shell}
    \item Firmware in \code{firmware/lab2/main.go}. App code in \code{lib/lab2/}.
    \item Leave a subscriber running in a terminal for the whole lab. It is your ground truth.
  \end{steps}
  \begin{hint}[Order of work]
    Sensor first, then the broker, then the app. Do not write Flutter code before a reading appears in \code{mosquitto\_sub}.
  \end{hint}
\end{frame}

\begin{frame}[embedded,eyebrow=Task I]{Read one sensor over I2C}
  \begin{columns}[T]
    \begin{column}{0.55\textwidth}
      \begin{steps}
        \item Pick a sensor with a driver in \code{tinygo.org/x/drivers}. The BME280 is cheap, common, and reads temperature.
        \item Wire it to the two I2C pins of your board, supply and ground.
        \item Configure the bus at 400 kHz, then the driver.
        \item Print one reading every two seconds, and read it back with \code{tinygo monitor}.
      \end{steps}
    \end{column}
    \begin{column}{0.41\textwidth}
      \kv{Done when}{}
      \begin{checklist}
        \item The serial line shows a new value every two seconds.
        \item The value changes when you breathe on the sensor.
        \item A read error prints a message instead of stopping the loop.
      \end{checklist}
    \end{column}
  \end{columns}
  \begin{takeaway}{A driver is the datasheet, already read.}
    Writing register numbers by hand is how a two-hour lab becomes a four-hour lab.
  \end{takeaway}
\end{frame}

\begin{frame}[embedded,fragile,eyebrow=Task I]{The shape of the sensor loop}
  \begin{columns}[T]
    \begin{column}{0.55\textwidth}
\begin{golang}
// tinygo.org/x/drivers/bme280
i2c := machine.I2C0
i2c.Configure(machine.I2CConfig{
    Frequency: 400 * machine.KHz})
s := bme280.New(i2c)  // 0x76 or 0x77
s.Configure()

for {
    mC, err := s.ReadTemperature()
    if err != nil { println(err.Error()) }
    println("c:", float32(mC)/1000)
    time.Sleep(2 * time.Second)
}
\end{golang}
    \end{column}
    \begin{column}{0.41\textwidth}
      \lead{Milli-degrees, not degrees.}{The drivers return thousandths of a degree Celsius. Divide once and keep one unit in your code.}
      \begin{warning}
        \code{Connected()} answers before the first read. Call it and print the answer: a sensor that never acknowledges is a wiring fault.
      \end{warning}
    \end{column}
  \end{columns}
\end{frame}

\begin{frame}[embedded,fragile,eyebrow=Task I]{No board, or no sensor}
  \vskip 1mm
\begin{golang}
// Same program, one line swapped. Everything downstream stays identical.
celsius := 20 + float32(seq%100)/10   // a ramp, not a constant
seq++
\end{golang}
  \vskip 2mm
  \begin{columns}[T]
    \begin{column}{0.48\textwidth}
      \kv{On Wokwi}{}
      \muted{Use the board Lab 1 named for you. If the part library carries an I2C sensor with a TinyGo driver, wire it in the diagram. If not, feed the loop the counter above and say so in your README.}
    \end{column}
    \begin{column}{0.48\textwidth}
      \kv{The broker, from Wokwi}{}
      \muted{The simulator runs in the browser, so a broker on your own machine is out of reach. Publish to a public test broker, on a topic prefix nobody else will use.}
    \end{column}
  \end{columns}
  \vskip 3mm
  \muted{A constant proves nothing. Publish a value that moves, so the number on the phone shows the chain is alive.}
\end{frame}

\begin{frame}[embedded,eyebrow=Task II]{Publish the reading over MQTT}
  \begin{columns}[T]
    \begin{column}{0.55\textwidth}
      \begin{steps}
        \item Join the network. On the ESP32-C3 and ESP32-S3 the radio comes from the espradio package, not from the standard library.
        \item Connect to the broker with a client id that carries your team name.
        \item Publish one JSON document on \code{course/<team>/reading} every two seconds, with QoS 1 and the retain flag off.
        \item Print the result of every step. A silent board tells you nothing.
      \end{steps}
    \end{column}
    \begin{column}{0.41\textwidth}
      \kv{Done when}{}
      \begin{checklist}
        \item \code{mosquitto\_sub} prints one line every two seconds.
        \item The payload parses as JSON and carries a sequence number.
        \item Killing the broker and starting it again brings the board back without a reflash.
      \end{checklist}
    \end{column}
  \end{columns}
\end{frame}

\begin{frame}[embedded,fragile,eyebrow=Task II]{Joining the network and the broker}
  \begin{columns}[T]
    \begin{column}{0.55\textwidth}
\begin{golang}
// drivers/netlink, drivers/net/mqtt
link, _ := netlink.New()
err := link.NetConnect(
    &netlink.ConnectParams{
        Ssid: ssid, Passphrase: pass,
    })
if err != nil {
    println("wifi:", err.Error())
}
opts := mqtt.NewClientOptions().
    AddBroker(broker).
    SetClientID("mec-" + team + "-board")
c := mqtt.NewClient(opts)
c.Connect().Wait()
\end{golang}
    \end{column}
    \begin{column}{0.41\textwidth}
      \lead{Treat this as a shape.}{The package paths for the radio move faster than the protocol does. Keep the sequence: join, connect, publish, check the error each time.}
      \begin{hint}[The address]
        \code{broker} is your laptop's address on the lab network, port 1883. Not \code{localhost}: that is the board.
      \end{hint}
    \end{column}
  \end{columns}
\end{frame}

\begin{frame}[embedded,fragile,eyebrow=Task II]{The topic and the payload}
  \vskip 1mm
\begin{golang}
topic := "course/" + team + "/reading"          // one team, one topic
p := fmt.Sprintf(`{"team":%q,"c":%.2f,"seq":%d}`, team, celsius, seq)
c.Publish(topic, 1, false, p)                   // QoS 1, retain off
time.Sleep(2 * time.Second)
\end{golang}
  \vskip 3mm
  \begin{columns}[T]
    \begin{column}{0.48\textwidth}
      \kv{Rules}{}
      \muted{Identity in the topic, data in the payload. Never publish the value as a topic level.\par\vskip 1.5mm
      The sequence number is yours, not the protocol's: it is the only way to see a gap.}
    \end{column}
    \begin{column}{0.48\textwidth}
      \kv{Why QoS 1}{}
      \muted{At QoS 1 the broker acknowledges each message, so a lost packet is resent. A duplicate is harmless here because every document carries its own sequence number.}
    \end{column}
  \end{columns}
\end{frame}

\begin{frame}[embedded,fragile,eyebrow=Task II]{A Pico without Wi-Fi: the laptop bridge}
  \vskip 1mm
\begin{shell}
# macOS: stty -f <port> 115200      Linux: stty -F <port> 115200
cat /dev/tty.usbmodem1101 | while IFS= read -r line; do
  mosquitto_pub -h localhost -q 1 -t course/t07/reading -m "$line"
done
\end{shell}
  \vskip 3mm
  \muted{The board prints the JSON document, one per line, exactly as a networked board would publish it. The laptop only forwards. Nothing else in the lab changes, which is the point: the phone cannot tell the two set-ups apart.}
  \vskip 2mm
  \begin{warning}
    \code{tinygo monitor} holds the serial port open. Stop it before you start the bridge.
  \end{warning}
\end{frame}

\begin{frame}[eyebrow=Task III]{Show the value live in the app}
  \begin{columns}[T]
    \begin{column}{0.55\textwidth}
      \begin{steps}
        \item Add a \code{ReadingBloc} with one event for an incoming document and a state holding the last value and its sequence number.
        \item Connect an \code{MqttServerClient} and subscribe to your team topic at QoS 1.
        \item Decode each message on the updates stream and add an event to the BLoC.
        \item Render the value in a \code{BlocBuilder}, with a visible waiting state before the first message.
        \item Reconnect with a growing delay when the connection drops.
      \end{steps}
    \end{column}
    \begin{column}{0.41\textwidth}
      \kv{Done when}{}
      \begin{checklist}
        \item The number on screen changes every two seconds.
        \item Unplugging the board stops the updates and the screen says so.
        \item Plugging it back in resumes updates with no restart of the app.
        \item The client is closed when the screen is disposed.
      \end{checklist}
    \end{column}
  \end{columns}
\end{frame}

\begin{frame}[fragile,eyebrow=Task III]{Connect, subscribe, decode}
\begin{dart}
final c = MqttServerClient('10.0.2.2', 'phone-$team')  // mqtt_client
  ..port = 1883
  ..keepAlivePeriod = 30
  ..onDisconnected = _scheduleReconnect;
await c.connect();
c.subscribe('course/$team/reading', MqttQos.atLeastOnce);
c.updates!.listen((events) {
  final m = events.first.payload as MqttPublishMessage;
  final text = MqttPublishPayload.bytesToStringAsString(m.payload.message);
  _bloc.add(ReadingReceived(jsonDecode(text)));
});
\end{dart}
  \vskip 2mm
  \muted{\code{10.0.2.2} is your machine from the Android emulator. A real phone uses the laptop's address.}
\end{frame}

\begin{frame}[fragile,eyebrow=Task III]{Render it, and come back after a drop}
\begin{dart}
BlocBuilder<ReadingBloc, ReadingState>(
  builder: (ctx, s) => Text(s.label),   // 'waiting' or '21.4 C'
)
void _scheduleReconnect() async {
  var delay = const Duration(seconds: 1);
  while (!_closed && !_connected) {
    await Future.delayed(delay);
    try { await c.connect(); } catch (_) {}
    delay *= 2;              // stop growing at a minute
  }
}
\end{dart}
  \vskip 1mm
  \muted{Two clients with the same id disconnect each other, which looks exactly like a broken network.}
\end{frame}

\begin{frame}[fragile,eyebrow=Acceptance]{What the grader will do}
  \vskip 1mm
\begin{shell}
mosquitto_sub -h localhost -q 1 -t 'course/+/reading' -v
\end{shell}
  \vskip 3mm
  \begin{columns}[T]
    \begin{column}{0.48\textwidth}
      \kv{The demo}{}
      \begin{checklist}
        \item The topic prints one document every two seconds.
        \item The app shows the same value, live.
        \item Pulling the board's cable stops the updates.
        \item Plugging it back in resumes them, with no app restart.
      \end{checklist}
    \end{column}
    \begin{column}{0.48\textwidth}
      \kv{In the repository}{}
      \muted{The firmware, the app code, and a README naming your team topic, your board or simulator, and the broker address you used.\par\vskip 1.5mm
      One commit per task. A screenshot of the app next to the subscriber terminal is the cheapest proof you can leave.}
    \end{column}
  \end{columns}
\end{frame}

\begin{frame}[eyebrow=Troubleshooting]{Five failures you will meet}
  \trouble{Connected() is false}{Wrong I2C address: scan the bus, 0x76 or 0x77.}
  \trouble{The bus scan prints nothing}{Nothing acknowledged. Check supply, ground, and the two I2C pins.}
  \trouble{Error: Connection refused}{Nothing listens on 1883. Start Mosquitto, bind the network address.}
  \trouble{Both clients reconnect in a loop}{One client id, two clients. Give each of them its own id.}
  \trouble{CLEARTEXT not permitted}{Android blocks a plain broker. Set usesCleartextTraffic true in the manifest.}
\end{frame}

\closingframe{Push before you leave.}{Branch lab-2 · one commit per task · your topic in the README}

\end{document}
